\section{Introduction}
\label{sec:intro}
Estimating distance, matching objects across views, and interpreting a change in viewpoint require relating images to a three-dimensional scene. SPAR evaluates these capabilities across spatial tasks and cognitive levels~\cite{spar}. We examine what geometric evidence a vision-language model receives and how the model uses that evidence during reasoning.

3DThinker investigates geometric imagination and continuous latent reasoning from limited views~\cite{thinker}. PIGeR starts from the pretrained 3DThinker-Mindcube checkpoint and retains the eight-step recurrent computation used in our configuration. We ask: \emph{can measured camera poses and explicit RGB-D representations make continuous spatial reasoning more effective?} The study changes the geometric evidence available to the model while preserving its latent recurrence.

PIGeR (Pose-Integrated Geometric Reasoning) integrates measured relative poses and frozen BIP3D RGB-D features into the existing visual prefix~\cite{bip3d}. We adapt CGMF from SpaceMind~\cite{spacemind} to use camera embeddings for geometric K/V conditioning and output gating. Fusion updates the output of the frozen Qwen visual merger before the embeddings replace image placeholders in the language input. Visual-token count, order, and dimensions remain unchanged. A residual scale initialized to zero preserves the original visual prefix at initialization. Language-side LoRA adapters~\cite{lora} and the fusion branch are then optimized jointly with the original assistant-sequence objective.

PIGeR adapts existing components to study explicit geometric evidence in continuous reasoning. It reuses the latent mechanism of 3DThinker, the BIP3D feature extractor, and the SpaceMind fusion design. The language decoder receives question text; CGMF and geometric sampling do not. Both visual branches receive the official image markers, which can provide indirect target cues. Geometry reaches the continuous latents through the enhanced visual prefix. The answer decoder also retains direct access to that prefix. A change in latent states therefore does not, by itself, show that the latent pathway contributes to an answer improvement.

The comparison with GENERIC separates camera-specific conditioning from access to the additional inputs. GENERIC receives the same RGB-D inputs, calibration, and camera embeddings as CGMF, with the same trainable parameter count and initialization. It replaces the camera-specific conditioning sources in the K/V bias and output gate while retaining camera embeddings in attention memory. The comparison therefore evaluates these conditioning choices jointly. B0 provides an RGB-only reference at a matched optimization budget. The fusion variants receive additional inputs and train additional parameters, which limits attribution of their differences from B0.

The completed 30K comparison does not support an overall benefit from camera-specific conditioning. Full scores are 48.530 for CGMF, 50.345 for GENERIC, and 49.489 for B0 (Table~\ref{tab:main}). The paired Overall interval for CGMF minus GENERIC is negative. Development-set donor substitutions provide evidence of sensitivity on pose tasks, but their Overall intervals include zero. These findings distinguish sensitivity to geometric inputs from an improvement over either comparison model.

We describe a token-preserving integration of measured pose and frozen RGB-D features, then compare camera-guided and generic fusion under matched inputs and trainable parameter counts. The checkpoint trajectory, task breakdown, and input interventions test different aspects of this design. Together, they show where camera-specific conditioning falls short in the evaluated configuration without implying that geometric inputs are generally ineffective.
